%% =====================================================================
%%  1 / 7   Introduction  引言                                   成员 1
%% =====================================================================
\section{Introduction}
\label{sec:intro}

Medical artificial intelligence has traditionally been conceived as an advisory technology that supports clinicians while leaving decisions and actions to them \citep{topol2019}. Recent large language models unsettle this boundary because they can plan, invoke external tools, and execute multi-step tasks in electronic health record environments. In this paper, a \emph{medical agent} means a system with those three capabilities; single-turn question-answering systems are excluded. On \textsc{MedAgentBench}, the strongest model completes 69.67\% of 300 physician-written clinical tasks \citep{jiang2025medagentbench}, while Agent Hospital simulates an entire care process involving LLM-powered patient, nurse, and doctor agents \citep{li2024agenthospital}. Medical AI is therefore beginning not only to advise but also to act.

This transition changes the category of error. Accuracy, sensitivity, and task success assess the quality of a judgement, but not the legitimacy of an action. A wrong recommendation can still be reviewed; an executed medication change has already altered the patient's record and potentially the patient's body. The levels-of-autonomy framework proposed by \citet{ye2025autonomy}, spanning informational tools (L0) to supervised agents (L3), recognizes that different actions carry different risks. However, it classifies what agents may do without fully specifying the conditions under which a clinician or institution may authorize them to do it.

Two largely separate literatures bear on this problem. Engineering research asks how agents can plan reliably, use tools correctly, and recover from failure. Human-factors research asks whether people remain able to supervise automation. \citet{bainbridge1983} showed that greater automation can make the remaining human tasks more critical and more difficult, and \citet{parasuraman2010} linked nominal supervision to predictable complacency and bias. In medicine, \citet{budzyn2025deskilling} found that endoscopists' adenoma detection rate in non-AI-assisted procedures fell from 28.4\% to 22.4\% after exposure to AI-assisted colonoscopy. This observational finding provides evidence that unaided competence may erode even when a human formally remains in the loop. In the literature reviewed here, capability and supervision are thus examined separately, while the authorization that connects them remains underspecified.

Against this background, this paper asks: \emph{under what conditions can clinical action be legitimately delegated to a medical AI agent?} We address this question through a purposive review of seventeen engineering, clinical, governance, and human-factors studies. The review identifies four connected themes---deskilling, the supervision illusion, accountability diffusion, and trust under uncertainty---and treats them as symptoms of an unspecified authorization boundary. We also specify, but do not yet execute, a $2\times2$ vignette experiment that crosses the system's role (advisor or agent) with disclosure of its AI identity. On this basis, we argue that the human--AI relationship should be framed not as vague \emph{human oversight} but as \emph{governed delegation}: a transfer of authority whose scope is stated in advance, which can be revoked during execution, whose actions remain auditable, and which preserves the human's unaided capacity. This framework is a conceptual specification to be tested, rather than an empirical effect already established.

Section~\ref{sec:litreview} establishes the gap, and Section~\ref{sec:method} describes the review and vignette protocol. Section~\ref{sec:results} reports the corpus and four themes; Section~\ref{sec:discussion} develops governed delegation and its implications. Section~\ref{sec:conclusion} concludes by asking not merely whether medical AI can act, but on what terms its actions can be made answerable.
